\section{Materiais e Métodos}
\label{sec:material_metodos}

A metodologia compreendeu a construção do nariz eletrônico, o delineamento
experimental, a aquisição dos sinais do sistema solo--planta, o
pré-processamento das séries temporais e a classificação das coletas a vácuo.
A validação agrupada por vaso foi usada para avaliar as coletas da base
completa; a partição histórica por coleta e a exclusão exploratória de
coletas foram examinadas separadamente.

\subsection{Sistema de nariz eletrônico}
\label{subsec:sistema_nariz}

O nariz eletrônico foi desenvolvido para registrar a resposta conjunta de
sensores de gases à amostra do espaço acima do solo. A matriz sensorial e as
grandezas auxiliares são descritas primeiro; em seguida, são apresentados a
estrutura de amostragem, o controle eletrônico, a interface e o armazenamento.

\subsubsection{Matriz de sensores de gases}
\label{subsubsec:matriz_sensores}

A matriz foi constituída por MQ-2, MQ-3, MQ-7, MQ-8, MQ-135 e MQ-138,
instalados na mesma câmara de vidro e expostos à amostra gasosa. Esses sensores
semicondutores de óxido metálico apresentam sensibilidades parcialmente
sobrepostas. Por isso, a classificação utilizou o padrão conjunto das seis
respostas, sem atribuir a um sensor isolado a identificação química de um
composto. As curvas de sensibilidade dos fabricantes foram empregadas apenas
para descrever as respostas nominais dos modelos MQ. As leituras registradas
foram sinais instrumentais, sem conversão para concentrações de compostos.

\begin{figure}[H]
    \centering
    \includegraphics[width=0.30\textwidth]{figures/matriz.png}
    \caption{Matriz dos seis sensores de gases na câmara sensorial.}
    \label{fig:matriz_sensores}
\end{figure}

\subsubsection{Estrutura física e sensores auxiliares}

O espaço de amostragem era delimitado por um tubo de PVC introduzido no
substrato do vaso. Uma extremidade do tubo possuía recortes triangulares,
que facilitavam a fixação. O sensor capacitivo de umidade do solo,
\textit{Capacitive Soil Moisture Sensor V2.0}, ficava na parte interna do
tubo, com a haste inserida no mesmo substrato amostrado. Sobre o tubo,
uma caixa de MDF sustentava a câmara sensorial de vidro, a placa dos
sensores MQ, os componentes eletrônicos e as conexões com as mangueiras.
O BMP280 registrava a temperatura e a pressão absoluta no equipamento.
Esse sensor não mede umidade relativa do ar; a leitura do sensor de solo
foi tratada como índice operacional, sem calibração em teor de água.
% CONFERIR: a qualificação menciona o BMP280 no interior da estrutura e
% também traz uma legenda que o situa dentro da câmara de gases. Confirmar
% sua posição física antes de especificá-la na versão de submissão.

\begin{figure}[H]
    \centering
    \includegraphics[width=0.30\textwidth]{figures/completo.jpeg}
    \caption{Estrutura externa do nariz eletrônico.}
    \label{fig:estrutura_nariz_eletronico}
\end{figure}

\subsubsection{Controle eletrônico e digitalização}

O ESP32 sequenciava as etapas da coleta, acionava a bomba de vácuo e as
válvulas por meio do módulo de relés e reunia as leituras dos sensores.
Os sinais analógicos da matriz MQ passavam por um conversor
analógico--digital externo antes de serem registrados pelo sistema.
% CONFERIR NO ESQUEMA OU FIRMWARE: modelo, resolução, tensão de referência,
% interface com o ESP32, número de canais e atribuição de cada MQ ao ADC.
% A qualificação anexada não identifica o conversor; não completar esses
% dados a partir de especificações genéricas do ESP32 ou dos sensores.

O sistema pneumático reunia a bomba, as válvulas, as mangueiras e a câmara
de vidro. O ESP32 acompanhava a pressão medida pelo BMP280 para interromper
a depressurização no valor de referência e, então, comandar a admissão da
amostra. O estado da bomba e a etapa em execução podiam ser acompanhados
pela interface de monitoramento.

\subsubsection{Interface, configuração e armazenamento}

Uma interface web acessada por Wi-Fi e navegador permitia operar o
equipamento por computador, tablet ou smartphone. A página inicial dava
acesso às páginas de controle, monitoramento e dados, além de apresentar
informações sobre o cartão SD e os recursos de memória. Na página de
controle, o operador definia os tempos de extração do ar, aquecimento,
obtenção da amostra, análise das respostas e desaturação, além da pressão
de referência e do nome do arquivo de saída; depois iniciava o ciclo.
Esses parâmetros eram configuráveis na interface. Os valores usados como
referência no protocolo analisado são apresentados na
Seção~\ref{subsec:protocolo_aquisicao}.

Durante a coleta, a página de monitoramento mostrava a etapa ativa, o
tempo decorrido, o estado da bomba, a temperatura, a pressão e informações
de armazenamento. Em cada leitura eram registrados o tempo decorrido, os
seis sinais MQ, o sinal do sensor de solo, a temperatura e a pressão.
O equipamento armazenava essas séries em arquivos TXT no cartão SD.
Ao final da coleta, a página de dados listava os arquivos e seus tamanhos
e permitia baixá-los para outro dispositivo ou excluí-los. Assim, o
computador podia ser usado para configuração, acompanhamento e recuperação
dos arquivos por meio do navegador.
% CONFERIR NO FIRMWARE: frequência efetiva de gravação no SD, formato e
% separador do TXT, marcação das etapas e tratamento de erro de escrita.

\subsubsection{Alimentação para uso em campo}

O equipamento podia receber energia de uma fonte externa de 19~V ligada à
rede elétrica ou de uma bateria de motocicleta de 12~V em locais sem rede.
Em ambos os casos, um regulador fornecia aproximadamente 5~V ao barramento
que alimentava os sensores MQ e os dispositivos de acionamento. O módulo
ESP32 recebia alimentação desse barramento e utilizava seu regulador
interno para operar a lógica em 3,3~V. A alimentação por bateria era uma
opção para locais sem rede elétrica. A capacidade da bateria, o consumo
medido e a duração de operação não foram quantificados neste estudo.
% CONFERIR: tensão e corrente efetivas sob carga, capacidade da bateria,
% consumo por etapa e autonomia medida antes de afirmar operação por um
% número determinado de coletas ou horas.

\subsection{Delineamento experimental e material vegetal}
\label{subsec:delineamento}

O experimento foi conduzido na Estação de Pesquisas Agronema, em Grandes
Rios, Paraná, Brasil, nas coordenadas $24{,}186^{\circ}$ S e
$51{,}399^{\circ}$ O. Foram utilizados dez vasos com plantas de soja,
distribuídos em dois grupos experimentais de cinco vasos cada: controle e
inoculado. Os vasos constituíram as unidades biológicas do experimento.
Leituras sucessivas e coletas repetidas do mesmo vaso não foram tratadas
como réplicas biológicas independentes.

As aquisições começaram aproximadamente 35 dias após a inoculação, em
11 de maio de 2026. A base utilizada para a análise com coleta a vácuo
contém 36 coletas realizadas em 11, 12, 13, 18 e 19 de maio de 2026,
correspondentes, aproximadamente, a 35, 36, 37, 42 e 43 dias após a
inoculação. Cada vaso foi amostrado em três ou quatro dessas ocasiões.
Das 36 coletas, 17 pertenciam ao grupo inoculado e 19 ao controle.
Houve irrigação entre as coletas de 13 e 18 de maio. O experimento não
estabeleceu níveis controlados de água no solo para esta comparação, nem
padronizou ou documentou a quantidade de água aplicada por vaso, o modo de
irrigação e o intervalo de estabilização entre a irrigação e cada coleta.
Essas informações não podem ser recuperadas a partir do índice do sensor.
Portanto, a comparação entre faixas do índice de solo foi tratada como
análise exploratória, e não como teste do efeito da irrigação ou da umidade.




\subsection{Vasos controle e vasos inoculados com \textit{Heterodera glycines}}
\label{subsec:classes_experimentais}

Conforme o delineamento experimental, os vasos 1 a 5 formaram o grupo
controle, denominado saudável, e não receberam nematoides. Os vasos 6 a 10
formaram o grupo inoculado e receberam aproximadamente 100 mil juvenis de
segundo estádio (J2) de
\textit{H. glycines} por vaso.

A Figura \ref{fig:vasos_experimento} apresenta os vasos
utilizados no experimento, divididos entre os grupos
controle e inoculado.

\begin{figure}[H]
    \centering
    \includegraphics[width=0.55\textwidth]{figures/vasos_experimento.png}
    \caption{Vasos de soja utilizados no experimento,
    divididos entre os grupos controle não inoculado
    e inoculado com \textit{H. glycines}.}
    \label{fig:vasos_experimento}

    \vspace{0.1cm}
    {\footnotesize Fonte: Elaborado pelos autores (2026).}
\end{figure}

Na base de classificação, a \textbf{classe 0} identifica o grupo
experimentalmente inoculado e a \textbf{classe 1} identifica o controle
não inoculado. Esses rótulos representam a alocação experimental, não a
confirmação nematológica individual de cada vaso. Neste artigo, o recall
da classe inoculada corresponde ao recall da classe 0.


\subsection{Protocolo de aquisição com pré-depressurização}
\label{subsec:protocolo_aquisicao}

O tubo de PVC foi acoplado ao vaso para delimitar o espaço de amostragem
(Figura~\ref{fig:vasos_experimentos}). Antes da entrada da amostra, a
bomba, controlada pelo ESP32, reduzia a pressão interna do sistema para
aproximadamente 70~kPa, enquanto a pressão externa estava próxima de
94~kPa. A finalidade dessa depressurização era criar uma diferença de
pressão que, na abertura da válvula, conduzisse inicialmente o ar do
espaço acima do solo até a câmara sensorial. O valor de 70~kPa era a
referência operacional configurada para encerrar essa preparação; não
representava a pressão mantida durante a leitura. Ao encerrar a
depressurização, o ESP32 desligava a bomba e abria a válvula de entrada.
A diferença de pressão conduzia inicialmente o ar do \textit{headspace}
do vaso pela tubulação até a câmara dos sensores, cuja pressão subia para
aproximadamente 94--95~kPa. Conforme o fluxograma de operação, a válvula
de entrada permanecia aberta durante a aquisição da amostra e a análise
das respostas, com a bomba de extração desligada. Depois, a válvula era
fechada e a bomba acionada para expulsar o ar da câmara na desaturação.
O protocolo não descreve um fluxo forçado contínuo durante a leitura.
% ARTUR: informar, se houver registro de ensaio ou especificação do projeto,
% por que a referência de 70 kPa foi escolhida. A documentação disponível
% explica sua função, mas não justifica esse valor numérico específico.
Uma execução
completa do ciclo, registrada sob um identificador próprio de vaso e dia, foi
definida como uma coleta. Em cada leitura, o ESP32 registrou o tempo decorrido,
os seis sinais MQ, a leitura do sensor de solo, a temperatura e a pressão.
Somente as coletas realizadas no protocolo de amostragem a vácuo integraram
a base analisada.

\begin{figure}[H]
    \centering
    \includegraphics[width=0.55\textwidth]{figures/vasos_experimentos.png}
    \caption{Protótipo posicionado sobre o vaso de soja durante a coleta a vácuo.}
    \label{fig:vasos_experimentos}
    \vspace{0.1cm}
    {\footnotesize Fonte: Elaborado pelos autores (2026).}
\end{figure}

A Figura~\ref{fig:ciclo_coleta} apresenta o ciclo de operação com
pré-depressurização. O termo ``coletas a vácuo'' designa aqui as
aquisições em que a câmara foi previamente depressurizada pela bomba e
recebeu a amostra após a abertura da válvula. Coletas de outros protocolos
não foram combinadas com elas nesta análise. A pressão registrada como
entrada dos classificadores corresponde às leituras retidas após o recorte
temporal e a seleção computacional descritos adiante, não ao valor-alvo de 70~kPa usado
na preparação da câmara.

\begin{figure}[H]
    \centering
    \includegraphics[width=\textwidth]{figures/fluxo_en.png}
    \caption{Fluxograma do ciclo de coleta com pré-depressurização
    e esquema do sistema de amostragem.}
    \label{fig:ciclo_coleta}
    \vspace{0.1cm}
    {\footnotesize Fonte: Elaborado pelos autores (2026).}
\end{figure}

O ciclo de referência, ilustrado na Figura~\ref{fig:ciclo_coleta},
compreendeu três minutos de aquecimento dos sensores; depressurização
até a referência de 70~kPa; abertura da válvula de entrada, com elevação
da pressão interna para aproximadamente 94--95~kPa; cinco minutos de
aquisição da amostra; cinco minutos de análise das respostas; e, após o
fechamento da válvula, três minutos de limpeza e desaturação com a bomba.
A válvula de entrada permanecia aberta durante os dois intervalos de cinco
minutos. Essas durações são os parâmetros do fluxograma, não marcas de
tempo verificadas individualmente nas 36 séries: os registros tiveram
durações diferentes. O controle da bomba e das válvulas foi realizado
pelo ESP32 durante o ciclo.
% CONFERIR A ORDEM NO FIRMWARE: a tabela da qualificação põe o aquecimento
% antes da extração; o parágrafo do protocolo sugere a ordem inversa.
% Verificar a configuração realmente usada nas 36 coletas.
A duração original de cada coleta foi definida como a diferença entre o
maior e o menor tempo decorrido registrados na planilha original da coleta;
essas durações foram resumidas separadamente por classe. Também foram
descritos, por classe, os valores de pressão antes e depois da filtragem.

A condição de solo seco ou úmido foi definida pelo índice operacional
\texttt{Soil\_indice\_0\_1}, obtido pela normalização mínimo--máximo da leitura
do sensor de umidade do solo na base consolidada, conforme a
Equação~\eqref{eq:indice_solo}:
\begin{equation}
    I_{\mathrm{solo}} =
    \frac{\mathrm{Soil}-\min(\mathrm{Soil})}
         {\max(\mathrm{Soil})-\min(\mathrm{Soil})},
    \label{eq:indice_solo}
\end{equation}
em que $\mathrm{Soil}$ é a leitura bruta do sensor de umidade do solo, e
$\min(\mathrm{Soil})$ e $\max(\mathrm{Soil})$ são, respectivamente, os valores
mínimo e máximo dessa leitura na base consolidada.
Esses extremos foram determinados antes das divisões de treino e teste;
portanto, as leituras posteriormente destinadas ao teste participaram da
definição da escala do índice. A escala não foi ajustada novamente dentro
de cada rodada de validação original. Uma reanálise LOSO com extremos
calculados somente no treino de cada rodada é descrita adiante.

Leituras com $I_{\mathrm{solo}}\leq 0{,}4$ foram denominadas \emph{solo seco};
leituras com $I_{\mathrm{solo}}>0{,}4$, \emph{solo úmido}. Esses nomes indicam
apenas faixas de um índice operacional, sem medição calibrada do teor de água
ou controle experimental da irrigação. No conjunto de
teste, todas as leituras de cada coleta ficaram do mesmo lado desse limiar,
permitindo classificar também a coleta inteira como seca ou úmida. O índice não
representa porcentagem calibrada de água no solo.

A Tabela~\ref{tab:solo_classe_dia} descreve a distribuição das 36
coletas por dia do mês, grupo experimental e faixa operacional do índice.
Para atribuir uma única faixa a cada coleta nessa tabela, utilizou-se a
mediana das suas leituras de $I_{\mathrm{solo}}$; duas coletas controle do
dia 13 apresentaram valores individuais ligeiramente acima de 0,4,
embora sua mediana permanecesse na faixa seca. Ambas pertenciam ao
treinamento da partição histórica por coleta; por isso, a afirmação anterior
de que cada coleta do teste ficou inteiramente em uma faixa continua válida.
As análises por leitura
continuaram usando o limiar em cada leitura.

\begin{table}[H]
    \centering
    \caption{Número de coletas por dia, classe e faixa operacional do índice de solo.}
    \label{tab:solo_classe_dia}
    \begin{tabular}{ccccc}
        \hline
        Dia do mês & Seco, classe 0 & Seco, classe 1 & Úmido, classe 0 & Úmido, classe 1 \\
        \hline
        11 & 2 & 5 & 0 & 0 \\
        12 & 5 & 1 & 0 & 0 \\
        13 & 0 & 4 & 0 & 0 \\
        18 & 0 & 0 & 5 & 4 \\
        19 & 0 & 0 & 5 & 5 \\
        \hline
    \end{tabular}
\end{table}

\subsection{Pré-processamento dos sinais e janelamento}
\label{subsec:preprocessamento}

Cada planilha da base original representou uma coleta. As leituras foram
associadas ao identificador da coleta, ao dia, ao vaso e à classe. Na
consolidação, foram removidos 15\% dos registros iniciais e 15\% dos
registros finais de cada série para reduzir a contribuição dos transientes
de início e fim do ciclo. Esse recorte proporcional não equivale a impor
os mesmos limites de tempo a todas as coletas: permaneceu a faixa central
de 70\% de cada série, que pode conter trechos fora da amostragem nominal
de 200 a 800~s. As colunas necessárias foram
convertidas para formato numérico, e registros sem valores válidos nessas
colunas foram descartados.

\subsubsection{Exclusão computacional de leituras instáveis pela pressão}

Após a coleta, as séries temporais passaram por uma seleção computacional
de leituras. Esse procedimento não corresponde a um filtro, válvula ou
dispositivo físico instalado no sistema pneumático. Seu objetivo foi
retirar trechos com mudanças bruscas de pressão e leituras distantes da
região de pressão observada após a admissão da amostra.

Primeiro, calculou-se
$\Delta P_t=|P_t-P_{t-1}|$ dentro de cada coleta. Quando
$\Delta P_t\geq 0{,}10$~kPa, foram excluídas a leitura do evento e até
30 leituras anteriores e 30 posteriores. Em seguida, na versão estrita
do procedimento, também foram removidas leituras fora da faixa definida
pela mediana global da pressão $\pm 0{,}50$~kPa, calculada após a exclusão
dos saltos na base
consolidada. Essa mediana foi 93,69~kPa, resultando na faixa de
93,19 a 94,19~kPa. Os valores de 93,69~kPa e da faixa correspondente
foram obtidos dos dados; não eram ajustes da bomba ou das válvulas. A
seleção reteve principalmente leituras próximas à pressão externa após
a entrada inicial da amostra, e não as leituras da depressurização até
cerca de 70~kPa. Após essas exclusões, a base de coletas a vácuo continha
63.492 leituras de 36 coletas. Não foi aplicada média móvel nem outra suavização temporal
aos sinais usados pelos classificadores.
% ARTUR/AUTORES: documentar a origem dos limiares de 0,10 kPa,
% 30 leituras e +/-0,50 kPa (ensaio preliminar, inspeção dos dados ou
% especificação prévia). Sem essa informação, não atribuir a eles uma
% justificativa física ou instrumental que não foi registrada.
A exclusão em janelas de até 30 leituras de cada lado do evento retirou
2.276 de 32.226 leituras da classe 0 (7,06\%) e 2.945 de 39.076 da
classe 1 (7,54\%). A restrição adicional pela faixa de pressão removeu
mais 906 leituras da classe 0 (2,81\% das 32.226 originais) e 1.683 da
classe 1 (4,31\% das 39.076 originais). Considerando as duas
etapas, foram retirados 9,87\% e 11,84\% das leituras originais dessas
classes, respectivamente.
A mediana de pressão também foi calculada antes das divisões de treino e
teste. A faixa central proporcional não corresponde necessariamente apenas
à fase nominal de exposição ao \textit{headspace}. A sensibilidade ao uso
da mediana global foi examinada na reanálise LOSO descrita adiante.

Na análise histórica com divisão por coleta, as entradas foram seis
sinais MQ corrigidos estatisticamente, o índice de solo, a temperatura em $^{\circ}$C e a
pressão em kPa. Para cada MQ, ajustou-se somente com as coletas de
treinamento um \texttt{HuberRegressor} que estimou a componente associada
ao índice de solo, à temperatura e à pressão. Essa componente foi
centralizada pela média das previsões no treinamento e subtraída do
sinal original, conforme a Equação~\eqref{eq:compensacao_mq}:
\begin{equation}
    \label{eq:compensacao_mq}
    MQ_{s,t}^{\mathrm{corr}} = MQ_{s,t} -
    \left(\widehat{MQ}_{s,t}^{\mathrm{amb}}-
    \overline{\widehat{MQ}_{s}^{\mathrm{amb}}}_{\mathrm{treino}}\right).
\end{equation}
Nessa equação, $s$ identifica o sensor MQ, $t$ identifica a leitura, e a
barra representa a média das previsões nas leituras de treinamento. Os
``MQ corrigidos'' são, portanto, sinais com compensação \emph{estatística}
da associação linear estimada com essas três variáveis. Eles não resultam
de calibração física para concentração de gases ou umidade relativa do ar.
Essa compensação pertence à análise histórica por coleta; as análises LOSO
descritas adiante utilizaram sinais MQ brutos. As seis entradas MQ, o
índice de solo, a temperatura e a pressão totalizaram nove canais por
leitura na análise histórica.

Para a TCN, as leituras foram ordenadas por coleta e tempo e segmentadas
em janelas de 128 leituras, com passo de 64 leituras e sobreposição nominal
de 50\%, sem atravessar limites entre coletas. Quando necessário, uma
janela adicional foi alinhada ao fim da coleta para incluir as últimas
leituras; sua sobreposição podia superar 50\%. Cada janela teve formato de
$9\times128$ (canais $\times$ leituras) na análise histórica. As entradas
da rede foram padronizadas com parâmetros estimados somente nas coletas de
treinamento. Na base processada, o intervalo mediano observado entre
leituras consecutivas da mesma coleta foi de aproximadamente 0,24~s;
como há lacunas de aquisição, 128 leituras não equivalem sempre à mesma
duração em segundos.


\subsection{Modelos de aprendizado de máquina}
\label{subsec:modelos}

Foram utilizados três classificadores em análises distintas:
\emph{Extra Trees} (ET), \emph{Random Forest} (RF) e uma rede
convolucional temporal (TCN). A validação principal por vaso empregou
ET; a partição histórica por coleta comparou ET e TCN; a análise
exploratória após cinco exclusões comparou os três modelos. O ET foi
implementado com
\texttt{ExtraTreesClassifier}, 900 árvores,
\texttt{max\_features="sqrt"}, mínimo de dez amostras por folha,
\texttt{bootstrap=False}, sem pesos de classe e semente 42.
O RF foi implementado com \texttt{RandomForestClassifier}, 300 árvores,
\texttt{class\_weight="balanced"} e semente 42, apenas na comparação
LOSO após exclusões. Os resultados históricos de 96,15\% no solo seco
e 74,46\% no solo úmido pertencem ao ET, não ao RF.

A TCN implementada neste trabalho usou três blocos convolucionais
temporais residuais, com 24, 24 e 32 canais, convoluções de núcleo 3 e
dilatações 1, 2 e 4. Cada bloco contém duas convoluções; assim, o campo
receptivo local corresponde a $1+4(1+2+4)=29$ leituras, cerca de 6,7~s
ao usar o intervalo mediano de 0,24~s. A agregação média global no tempo
combina as saídas locais de toda a janela de 128 leituras; depois vêm
\emph{dropout} de 0,20 e camada linear de duas classes. O
treinamento utilizou 60 épocas, lotes de 32 janelas, otimizador AdamW com
taxa de aprendizado $10^{-3}$ e \emph{weight decay} de $10^{-4}$, e
entropia cruzada com pesos inversamente proporcionais à frequência das
classes nas janelas de treinamento. A semente aleatória foi 42. A rede
produziu probabilidades por janela; para avaliar cada leitura, foi usada
a média das probabilidades das janelas que a cobriam, com limiar 0,5 para
a classe 1. A arquitetura implementada não incluiu mecanismo de atenção
nem método de explicabilidade integrado; por isso, é denominada TCN, e
não X-TCN. As 60 épocas e os demais hiperparâmetros foram mantidos fixos
nesta comparação, sem seleção por validação interna ou parada antecipada.

Na análise histórica, ET e TCN foram ajustados uma vez com as
coletas de treinamento. A comparação exploratória entre as faixas do índice
de solo foi feita após a predição, em dois recortes do mesmo conjunto de
teste, não em dois treinamentos independentes. Nas análises
LOSO, os modelos foram reajustados em cada rodada, conforme descrito adiante.

\subsection{Análise histórica com separação por coleta}
\label{subsec:treinamento_teste}

A divisão foi feita por \emph{coleta}, dentro de cada classe, com
proporção nominal de 70\% para treinamento e 30\% para teste e semente 42.
Todas as leituras e janelas de uma coleta permaneceram em um único conjunto.
Das 36 coletas, 24 foram destinadas ao treinamento (41.377 leituras: 11
coletas da classe 0 e 13 da classe 1) e 12 ao teste (22.115 leituras:
seis coletas de cada classe).

Essa separação evita que janelas da mesma aquisição apareçam nos dois
conjuntos. Entretanto, oito dos dez vasos aparecem em ambos porque foram
medidos em dias distintos. Assim, o teste avalia novas \emph{coletas}, não vasos
biologicamente independentes, e as métricas não demonstram generalização
para plantas ou vasos nunca vistos pelo modelo. Essa partição foi mantida
como comparação histórica; a validação agrupada por vaso foi usada como
referência principal para as análises com sinais MQ brutos.

No teste, o recorte seco incluiu 16.692 leituras de oito coletas: 6.507
leituras da classe 0 e 10.185 da classe 1. O recorte úmido incluiu 5.423
leituras de quatro coletas: 4.152 da classe 0 e 1.271 da classe 1. Essa
distribuição desigual motivou o relato da acurácia balanceada junto da
acurácia convencional. Como análise adicional da partição histórica,
foram usadas previsões do ET já obtidas no teste para formar um subconjunto
com igual número de leituras das duas classes em cada faixa: 6.507 de
cada classe na faixa seca e 1.271 de cada classe na faixa úmida. Todas
as leituras da classe minoritária foram mantidas e a classe majoritária
foi subamostrada com semente 42; as leituras excedentes foram examinadas
separadamente. Essa subamostragem ocorreu somente depois da predição,
sem retreinamento, e foi tratada como análise de sensibilidade a um
sorteio específico, não como resultado principal.

\subsection{Métricas da análise histórica}
\label{subsec:metricas}

Na partição histórica foram calculadas acurácia, acurácia balanceada,
F1 macro e matrizes de confusão no conjunto de teste. A acurácia é a
fração de leituras
classificadas corretamente; a acurácia balanceada é a média dos recalls
das classes 0 e 1; e a F1 macro é a média simples das medidas F1 das duas
classes. O recall da classe inoculada corresponde à sensibilidade
para a classe 0, ou seja, à fração de leituras de vasos inoculados
corretamente identificadas. As métricas foram calculadas separadamente
para os recortes seco e úmido do mesmo teste. A avaliação por
subamostragem balanceada foi descrita separadamente da análise do teste
completo.

\subsection{Sequência das análises}
\label{subsec:sequencia_analises}

A Tabela~\ref{tab:cronologia_analises} distingue as bases e as decisões
tomadas em cada etapa. A busca de subconjuntos e a inspeção dos erros
ocorreram antes da análise pós-exclusão; os controles de entradas fixas
foram acrescentados posteriormente para examinar a dependência dos
sinais ambientais.

\begin{table}[H]
    \centering
    \footnotesize
    \caption{Sequência das análises, base utilizada e decisão correspondente.}
    \label{tab:cronologia_analises}
    \begin{tabular}{p{0.07\textwidth}p{0.20\textwidth}p{0.61\textwidth}}
        \hline
        Etapa & Base & Objetivo e decisão \\
        \hline
        1 & 36 coletas & Recorte temporal e exclusão computacional de leituras instáveis pela pressão; definição da base comum. \\
        2 & 36 coletas & Partição histórica por coleta; comparação dos modelos com MQ corrigidos e avaliação adicional por condição do solo. \\
        3 & 36 coletas & Busca LOSO de 63 subconjuntos MQ com solo, temperatura e pressão fixos; seleção do conjunto MQ2, MQ3, MQ8, MQ135 e MQ138 e obtenção das previsões fora de cada vaso. \\
        4 & 36 coletas & Auditoria de saltos e de erros por coleta; definição exploratória das cinco coletas excluídas apenas da análise de sensibilidade. \\
        5 & 31 coletas & LOSO após exclusões; comparação de ET, Random Forest e TCN em entradas fixas; buscas ET de 63 subconjuntos MQ em cada um de dois cenários e das 511 combinações de entradas. \\
        6 & 36 coletas & Controles LOSO com todas as nove entradas, somente MQ, somente variáveis ambientais e regressão logística, sem seleção de subconjunto nessa etapa. \\
        7 & 36 coletas ou recortes por dia/tempo & Sensibilidade por dia, período e fase temporal; comparação pareada das configurações fixas. \\
        \hline
    \end{tabular}
\end{table}


\subsection{Auditoria exploratória das coletas}
\label{subsec:auditoria_coletas_loso}

Além da seleção de leituras pela pressão descrita acima, examinou-se a estabilidade dos
seis sinais MQ em cada coleta. Para o sensor $s$, a coleta $c$ e a leitura
$t>1$, calcularam-se a variação absoluta entre leituras consecutivas e o
limiar de cada sensor conforme a Equação~\eqref{eq:salto_mq_loso}:
\begin{equation}
    d_{s,c,t}=\left|x_{s,c,t}-x_{s,c,t-1}\right|,
    \qquad
    q_s=Q_{0{,}995}\!\left(\{d_{s,c,t}:c,\ t>1\}\right),
    \label{eq:salto_mq_loso}
\end{equation}
em que $x_{s,c,t}$ é o sinal bruto do MQ e $Q_{0{,}995}$ é o percentil
99,5 das variações reunidas de todas as coletas para aquele sensor.
Os limiares $q_s$ foram estimados globalmente antes das partições LOSO,
incluindo as coletas que depois apareceriam como teste em cada rodada.
As diferenças foram calculadas dentro de cada coleta, sem conectar o fim
de uma coleta ao início de outra. Definiu-se uma leitura com salto
multissensor quando pelo menos dois dos seis MQ superaram seus respectivos
limiares. O indicador de salto e a fração de leituras sinalizadas na coleta
foram calculados pela Equação~\eqref{eq:ruido_multissensor_loso}:
\begin{equation}
    J_{c,t}=\mathbf{1}\!\left\{
      \sum_{s\in\mathcal{S}_{\mathrm{MQ}}}
      \mathbf{1}\{d_{s,c,t}\geq q_s\}\geq 2
    \right\},
    \qquad
    R_c=\frac{1}{n_c}\sum_{t=1}^{n_c}J_{c,t},
    \label{eq:ruido_multissensor_loso}
\end{equation}
em que $\mathcal{S}_{\mathrm{MQ}}$ é o conjunto dos seis sensores MQ e
$n_c$ é o número de leituras da coleta. A regra exploratória de
qualidade sinalizou coletas com $R_c\geq 0{,}05$. A avaliação visual das
séries foi usada para contextualizar os saltos. O limiar foi aplicado à
coleta individual, de forma que as demais coletas do mesmo vaso permaneceram
elegíveis. Os percentis e a regra de exclusão foram definidos após a
inspeção da base completa; sua utilização constitui análise de
sensibilidade, e não um protocolo de qualidade estabelecido antes da
aquisição.
Saltos abruptos podem corresponder a respostas reais do sistema amostrado;
essa regra não identifica, por si só, falha instrumental.

Antes de qualquer exclusão de coletas, os 63 subconjuntos não vazios dos
seis MQ foram avaliados com as três variáveis ambientais fixas na base
completa, mediante dez rodadas LOSO. O conjunto MQ2, MQ3, MQ8, MQ135 e
MQ138 foi escolhido nessa busca pelo desempenho por leitura; o MQ7 não
foi retirado por motivo instrumental, apenas não integrou o subconjunto
selecionado. As previsões desse conjunto no LOSO foram obtidas no vaso
de teste de cada rodada, de modo que todas as 36 coletas receberam
previsões fora do treinamento. Portanto, os erros descritos a seguir
não vieram apenas das 12 coletas de teste da partição histórica nem de
previsões feitas nas amostras de treino. A fração de erros de cada coleta
foi calculada pela
Equação~\eqref{eq:erro_coleta_loso}:
\begin{equation}
    E_c=\frac{1}{n_c}\sum_{t=1}^{n_c}
    \mathbf{1}\{\widehat y_{c,t}\neq y_c\},
    \label{eq:erro_coleta_loso}
\end{equation}
em que $y_c$ é a classe experimental da coleta. Foram examinadas as
coletas com $E_c>0{,}5$. Na análise de sensibilidade com exclusões,
retiraram-se as quatro coletas identificadas por esse critério e uma
coleta adicional sinalizada pela auditoria multissensor; uma das quatro
também satisfazia o critério de ruído. A seleção por $E_c$ utilizou os
erros do próprio classificador na mesma base, e não demonstra defeito de
aquisição nas coletas selecionadas. Por isso, todas as comparações após
essas exclusões são exploratórias e condicionadas a essa seleção.

A Tabela~\ref{tab:coletas_excluidas_sensibilidade} identifica as coletas
retiradas somente da análise de sensibilidade. O vaso 10 inoculado aparece
como ``vaso 1'' em uma planilha derivada; sua identificação física foi
recuperada da planilha de origem.
\begin{table}[H]
    \centering
    \footnotesize
    \caption{Coletas retiradas da análise de sensibilidade, com o motivo da seleção.}
    \label{tab:coletas_excluidas_sensibilidade}
    \begin{tabular}{cccp{0.42\textwidth}}
        \hline
        Dia de maio & Classe & Vaso físico & Motivo \\
        \hline
        12 & 1, controle & 2 & Maioria das leituras incorreta no LOSO \\
        19 & 0, inoculado & 10 & Ruído multissensor \\
        19 & 0, inoculado & 9 & Maioria das leituras incorreta no LOSO \\
        19 & 1, controle & 1 & Maioria incorreta e ruído multissensor \\
        19 & 1, controle & 2 & Maioria das leituras incorreta no LOSO \\
        \hline
    \end{tabular}
\end{table}

\subsection{Validação principal com coletas agrupadas por vaso}
\label{subsec:loso_complementar}

Para avaliar coletas de vasos não vistos durante o ajuste, aplicou-se a
validação \textit{leave-one-vase-out} (LOSO) inicialmente à base completa,
sem excluir coletas. Todas as coletas de um mesmo vaso foram agrupadas:
em cada rodada, as coletas de um vaso constituíram o teste e as dos outros
nove constituíram o treinamento. Na base derivada, um vaso da classe
inoculada aparece com código numérico 1, embora a planilha de origem o
identifique como vaso 10. Por isso, a chave técnica de agrupamento combinou
classe e código de vaso, evitando colisão com o vaso 1 controle e mantendo
dez grupos distintos. Para cada vaso $v$, a separação das coletas de
treinamento e de teste seguiu a
Equação~\eqref{eq:particao_loso}:
\begin{equation}
    \begin{aligned}
        \mathcal{C}_v&=\{c\in\mathcal{C}:v(c)=v\},\\
        \mathcal{D}^{\mathrm{teste}}_v
        &=\bigcup_{c\in\mathcal{C}_v}\mathcal{D}_c,\\
        \mathcal{D}^{\mathrm{treino}}_v
        &=\bigcup_{c\in\mathcal{C}\setminus\mathcal{C}_v}\mathcal{D}_c,
    \end{aligned}
    \label{eq:particao_loso}
\end{equation}
em que $\mathcal{C}$ representa as coletas incluídas no cenário analisado,
$v(c)$ identifica o vaso da coleta $c$ e $\mathcal{D}_c$ reúne suas
leituras e classes. A base completa e a base da análise de sensibilidade
foram avaliadas separadamente com essa mesma regra. No segundo cenário,
retiraram-se apenas as cinco coletas descritas acima, mantendo os dez
vasos e suas demais coletas. Em nenhuma rodada coletas do vaso de teste
entraram no treinamento; cada leitura incluída foi testada uma vez.

Na base completa, o ET foi avaliado com MQ2, MQ3, MQ8, MQ135 e MQ138 mais
o índice de solo, a temperatura e a pressão. Essa configuração havia sido
escolhida entre subconjuntos de MQ avaliados nos mesmos vasos; o LOSO da
base completa é uma avaliação interna, não uma validação externa da
escolha das entradas. Os sinais MQ foram usados em sua forma bruta, sem
a compensação da Equação~\eqref{eq:compensacao_mq}. O ET foi reajustado em
cada rodada com 900 árvores,
\texttt{max\_features="sqrt"}, \texttt{min\_samples\_leaf=10},
\texttt{bootstrap=False} e semente 42. A previsão por leitura foi a classe
com probabilidade estimada de classe 1 igual ou superior a 0,5.

Para distinguir a contribuição das entradas sem selecionar previamente
um subconjunto, o mesmo LOSO na base completa foi repetido com três
configurações fixas de ET: os seis MQ mais solo, temperatura e pressão;
somente os seis MQ; e somente solo, temperatura e pressão. Também se
avaliou, com as nove entradas, uma regressão logística como modelo
simples de referência. Nessa regressão, a padronização das entradas foi
ajustada somente nas coletas de treinamento de cada rodada; usaram-se
\texttt{LogisticRegression}, \texttt{max\_iter=2000}, sem pesos de classe,
e semente 42. Esses controles compartilham as mesmas dez divisões por
vaso e os mesmos pré-processamentos globais já descritos; não constituem
uma validação externa do protocolo experimental.

Para examinar o efeito dos parâmetros globais de pré-processamento,
repetiu-se o LOSO com as 36 coletas antes da faixa estrita de pressão,
mas após o recorte temporal e o filtro fixo de saltos; nessa etapa havia
66.081 leituras. Em cada rodada, a mediana de pressão foi estimada
somente nas leituras dos nove vasos de treinamento. No cenário estrito,
aplicou-se a faixa mediana do treino $\pm 0{,}50$~kPa tanto ao treino
quanto ao vaso reservado para teste. No cenário de comparação, essa
restrição adicional não foi aplicada. Em ambos, o mínimo e o máximo
da leitura bruta de solo foram calculados somente no treino após a
eventual filtragem e usados para transformar treino e teste; valores
do teste fora da faixa $[0,1]$ não foram truncados. A exclusão de saltos,
definido por uma regra fixa dentro de cada coleta, foi mantido nos dois
cenários. Reajustou-se o mesmo ET de 900 árvores em cada rodada, com
quatro conjuntos de entradas: nove entradas, somente as três ambientais,
somente os seis MQ e os cinco MQ selecionados anteriormente mais as
três ambientais. O último conjunto permanece exploratório, pois sua
seleção não foi aninhada. As métricas por leitura, coleta e vaso foram
calculadas separadamente para cada cenário, informando-se os respectivos
denominadores porque a faixa estrita altera as leituras elegíveis.

Para avaliar a dependência das datas de aquisição, realizou-se ainda
uma análise \textit{leave-one-date-out} nas cinco datas de maio de 2026:
11, 12, 13, 18 e 19. Em cada rodada, todas as coletas de uma data
constituíram o teste, e as coletas das outras quatro datas constituíram
o treinamento. Reajustou-se o mesmo ET com as nove entradas fixas e,
como controles, somente as três entradas ambientais ou somente os seis
MQ. O mínimo e o máximo do solo e a mediana de pressão foram estimados
exclusivamente nas datas de treinamento de cada rodada. Repetiu-se a
análise com e sem a faixa estrita definida pela mediana do treino
$\pm 0{,}50$~kPa, mantendo em ambos os cenários o filtro fixo de saltos.
As previsões foram resumidas por leitura e por coleta, inclusive por
data reservada. Como um mesmo vaso pôde aparecer em diferentes datas,
essa análise examina transferência temporal e não substitui a validação
por vaso para unidades biológicas não vistas.

Na base após exclusão das cinco coletas, ET, \textit{Random Forest}
(\texttt{RandomForestClassifier}) e TCN foram reajustados em cada rodada
LOSO em duas configurações fixas: MQ2, MQ3, MQ8, MQ135 e MQ138, isoladamente ou
acompanhados das três variáveis ambientais. O Random Forest usou 300
árvores, \texttt{class\_weight="balanced"} e semente 42. Na rede temporal,
a padronização das entradas foi estimada apenas no treinamento de cada
rodada. As janelas mantiveram 128 leituras e passo de 64, com cinco ou oito
canais conforme a configuração; nenhuma janela atravessou o limite de
uma coleta. Para uma leitura $i$ coberta pelo conjunto $\mathcal{W}(i)$
de janelas, a probabilidade da classe 1 e a previsão final foram obtidas
pela Equação~\eqref{eq:agregacao_janelas_loso}:
\begin{equation}
    \widehat p_i(1)=
    \frac{1}{|\mathcal{W}(i)|}
    \sum_{w\in\mathcal{W}(i)}\widehat p_w(1),
    \qquad
    \widehat y_i=\mathbf{1}\{\widehat p_i(1)\geq 0{,}5\}.
    \label{eq:agregacao_janelas_loso}
\end{equation}
Os modelos foram comparados por leitura após reunir as previsões de todas
as rodadas. Essa avaliação não corresponde à partição por coleta de 70\%/30\%
descrita anteriormente. No treinamento e no teste LOSO, não houve
subamostragem, sobreamostragem nem geração de leituras sintéticas para
igualar o tamanho das classes. A subamostragem da partição histórica
ocorreu apenas depois das previsões. O ET
foi treinado sem pesos de classe; o Random Forest usou
\texttt{class\_weight="balanced"}, e a TCN usou pesos inversamente
proporcionais à frequência das classes nas janelas de treinamento.
Portanto, os três modelos não receberam o mesmo tratamento do
desbalanceamento. A acurácia balanceada, o F1 macro e os recalls por
classe foram relatados junto da acurácia simples para avaliar as duas
classes separadamente; essas métricas não modificam o treinamento do ET.

\subsection{Busca combinatória das entradas}
\label{subsec:busca_combinatoria_loso}

A busca exaustiva de entradas foi realizada somente com ET, na base após
as cinco exclusões, usando a mesma divisão LOSO por vaso. Ela não foi
aplicada ao Random Forest nem à TCN. Primeiro, examinaram-se todos os
subconjuntos não vazios dos seis sensores MQ, em dois cenários: (i) apenas
os MQ do subconjunto e (ii) os mesmos MQ acompanhados de índice de solo,
temperatura e pressão, sempre com as três variáveis ambientais juntas.
O número de subconjuntos MQ por cenário foi calculado pela
Equação~\eqref{eq:combinacoes_mq_loso}:
\begin{equation}
    N_{\mathrm{MQ}}=2^{6}-1=63.
    \label{eq:combinacoes_mq_loso}
\end{equation}
Assim, houve 63 configurações somente MQ e 63 com ambiente fixo. Cada
configuração foi avaliada nas mesmas dez rodadas LOSO e ordenada pela
acurácia por leitura reunida entre elas. Em seguida, a busca foi ampliada
para permitir a entrada separada de solo, temperatura e pressão. As duas
buscas de 63 configurações são recortes do espaço de busca ampliado.
As nove variáveis candidatas estão definidas na
Equação~\eqref{eq:conjunto_entradas_loso}:
\begin{equation}
    \mathcal{F}=\{\mathrm{MQ2},\mathrm{MQ3},\mathrm{MQ7},\mathrm{MQ8},
    \mathrm{MQ135},\mathrm{MQ138},I_{\mathrm{solo}},T,P\}.
    \label{eq:conjunto_entradas_loso}
\end{equation}
Foram avaliados todos os subconjuntos não vazios $A\subseteq\mathcal{F}$.
Cada variável pode estar presente ou ausente. Das $2^9$ possibilidades,
uma corresponde ao conjunto vazio, sem entradas para o classificador;
por isso, ela foi excluída.
O número de configurações foi calculado pela
Equação~\eqref{eq:combinacoes_loso}:
\begin{equation}
    \left|\{A\subseteq\mathcal{F}:A\neq\emptyset\}\right|
    =2^{|\mathcal{F}|}-1=2^9-1=511.
    \label{eq:combinacoes_loso}
\end{equation}
Solo, temperatura e pressão variaram separadamente nessa busca,
inclusive nas configurações sem sensores MQ. O BMP280 fornece
temperatura e pressão como duas variáveis, embora seja um único componente
físico; portanto, as 511 configurações correspondem a combinações de
variáveis de entrada, não de componentes. Se temperatura e pressão forem
tratadas como um bloco indivisível, há oito componentes selecionáveis:
seis sensores MQ, o sensor de solo e o BMP280. Cada componente pode estar
presente ou ausente. Excluindo a combinação sem nenhum componente, o
número de configurações físicas é dado pela
Equação~\eqref{eq:combinacoes_componentes_loso}:
\begin{equation}
    N_{\mathrm{componentes}}=2^8-1=255.
    \label{eq:combinacoes_componentes_loso}
\end{equation}
Esse total é uma contagem teórica de opções de hardware; a busca
computacional avaliou os 511 subconjuntos de variáveis.

Na busca das variáveis de entrada, cada subconjunto não vazio
$A\subseteq\mathcal{F}$ foi avaliado uma única vez. A ordem em que os
sensores aparecem não definiu novas configurações: conjuntos com as
mesmas variáveis foram considerados a mesma combinação, e as colunas
foram fornecidas ao classificador em uma ordem fixa. Para cada
combinação $A$, o Extra Trees foi reajustado nas dez rodadas LOSO,
usando somente as variáveis pertencentes a $A$. A acurácia dessa
combinação foi calculada após reunir as previsões por leitura de
todas as rodadas, conforme a Equação~\eqref{eq:acuracia_combinacao_loso}:
\begin{equation}
    \mathrm{Acc}(A)=
    \frac{
        \displaystyle\sum_{v\in\mathcal{V}}
        \sum_{i\in\mathcal{I}_v}
        \mathbf{1}\{\widehat{y}_{i}^{(A,v)}=y_i\}
    }{
        \displaystyle\sum_{v\in\mathcal{V}}|\mathcal{I}_v|
    },
    \label{eq:acuracia_combinacao_loso}
\end{equation}
em que $\mathcal{V}$ é o conjunto de vasos, $\mathcal{I}_v$ contém os
índices das leituras das coletas restantes do vaso reservado para
teste na rodada $v$, $y_i$ é a classe experimental e
$\widehat{y}_{i}^{(A,v)}$ é a previsão obtida com a combinação $A$
após treinar o modelo com as coletas dos outros nove vasos. Assim,
a mesma equação foi aplicada às 511 combinações; a ordem das colunas
não constituiu uma configuração distinta.

As combinações foram ordenadas pela acurácia por leitura. Em caso de
empate, preferiu-se a configuração com menos variáveis; persistindo
o empate, a que continha mais sensores MQ e, em seguida, a ordem
alfabética dos nomes das entradas. Como análise complementar do mesmo
ranking, a acurácia balanceada de cada combinação foi calculada a partir
dos recalls das duas classes, pela Equação~\eqref{eq:metricas_resumo_loso},
sem treinar novamente os modelos. A escolha original da configuração
utilizou a acurácia por leitura; o ranking balanceado foi examinado como
análise de sensibilidade. A escolha da configuração
e sua avaliação ocorreram nos mesmos vasos, sem validação aninhada
nem conjunto biológico externo; portanto, a ordenação das
configurações é exploratória.

\subsection{Métricas da validação por vaso}
\label{subsec:metricas_loso_complementar}

Com as previsões das rodadas LOSO reunidas, a acurácia por leitura foi
calculada conforme a Equação~\eqref{eq:acuracia_loso}:
\begin{equation}
    \mathrm{Acc}=\frac{1}{N}
    \sum_{i=1}^{N}\mathbf{1}\{\widehat y_i=y_i\},
    \label{eq:acuracia_loso}
\end{equation}
em que $N$ é o total de leituras avaliadas. Para cada classe
$k\in\{0,1\}$, o recall, a precisão e a medida F1 foram calculados
pela Equação~\eqref{eq:metricas_classe_loso}:
\begin{equation}
    \mathrm{Recall}_k=\frac{TP_k}{TP_k+FN_k},\qquad
    \mathrm{Prec}_k=\frac{TP_k}{TP_k+FP_k},\qquad
    F1_k=\frac{2\,\mathrm{Prec}_k\,\mathrm{Recall}_k}
    {\mathrm{Prec}_k+\mathrm{Recall}_k},
    \label{eq:metricas_classe_loso}
\end{equation}
em que $TP_k$, $FN_k$ e $FP_k$ são, respectivamente, verdadeiros positivos,
falsos negativos e falsos positivos ao tomar $k$ como classe positiva.
A acurácia balanceada e a média macro de F1 foram calculadas conforme a
Equação~\eqref{eq:metricas_resumo_loso}:
\begin{equation}
    \mathrm{Acc}_{\mathrm{bal}}=
    \frac{\mathrm{Recall}_0+\mathrm{Recall}_1}{2},
    \qquad
    F1_{\mathrm{macro}}=\frac{F1_0+F1_1}{2}.
    \label{eq:metricas_resumo_loso}
\end{equation}
O recall da classe 0 descreve a identificação de leituras do grupo
inoculado; nas métricas por leitura, os denominadores são leituras, não
vasos. As leituras sucessivas de uma mesma coleta e as coletas repetidas do mesmo
vaso não constituem réplicas biológicas independentes. Em particular,
a acurácia reunida atribui mais peso a coletas com maior número de leituras.

Para o ET com as nove entradas fixas na base completa, as previsões LOSO também
foram agregadas em decisões por coleta e por vaso. Calculou-se a média
das probabilidades das leituras de cada coleta; em seguida, as probabilidades
das coletas de um vaso receberam o mesmo peso para formar sua decisão,
conforme a Equação~\eqref{eq:agregacao_coleta_vaso}:
\begin{equation}
    \widehat p_c(1)=\frac{1}{n_c}\sum_{i\in\mathcal I_c}\widehat p_i(1),
    \qquad
    \widehat p_v(1)=\frac{1}{|\mathcal C_v|}
    \sum_{c\in\mathcal C_v}\widehat p_c(1).
    \label{eq:agregacao_coleta_vaso}
\end{equation}
em que $\mathcal I_c$ contém as leituras da coleta $c$ e $n_c$ é o seu
número de leituras. As decisões de coleta e vaso utilizaram limiar de
0,5 para a classe 1.
As mesmas métricas foram calculadas separadamente nas três unidades:
leitura, coleta e vaso. Na unidade vaso, a proporção de acertos foi
expressa como $k/10$ e acompanhada por intervalo binomial exato de
95\% pelo método de Clopper--Pearson. Como medida exploratória adicional,
para a acurácia balanceada foi calculado um intervalo por 10.000
reamostragens estratificadas com reposição de
cinco vasos por classe; cada vaso levou consigo todas as suas previsões
fora de treino. Usaram-se os percentis 2,5 e 97,5 das estimativas
reamostradas, com semente 42. Esses intervalos são condicionais ao modelo
e às entradas já escolhidas, sem incorporar a incerteza da seleção da
configuração entre os subconjuntos avaliados. Com apenas dez vasos, o
intervalo por reamostragem foi tratado como descritivo.

Para os controles fixos de ET e a regressão logística na base completa,
assim como para a comparação ET, Random Forest e TCN na base com cinco
coletas excluídas, calculou-se a área sob a curva ROC (AUC) por coleta a
partir das probabilidades médias $\widehat p_c(1)$ da
Equação~\eqref{eq:agregacao_coleta_vaso}. Essa métrica não depende do
limiar de classificação de 0,5. Como as probabilidades de cada vaso
provêm de um modelo ajustado em uma rodada diferente, a AUC reunida
entre rodadas foi tratada como descritiva. Não foi realizado teste pareado
por vaso para afirmar superioridade entre ET, Random Forest e TCN.

\subsection{Análises de sensibilidade para dia, período e fase da aquisição}
\label{subsec:sensibilidade_dia_periodo}

Para examinar a associação entre dia e classe, calculou-se primeiro um
referencial descritivo que atribuiu a cada coleta a classe majoritária
observada naquele dia em todas as 36 coletas. Como esse referencial usa
os próprios rótulos das coletas avaliadas, não constitui validação
fora de treino. Calculou-se também um referencial por dia fora de vaso:
para cada vaso reservado no LOSO, a classe atribuída às suas coletas foi
a majoritária entre as coletas dos outros nove vasos no mesmo dia.
Havendo empate nesse dia, utilizou-se a classe majoritária entre todas
as coletas dos nove vasos de treinamento. Cada coleta teve peso igual
nessa contagem, independentemente do seu número de leituras.

O ET com entradas previamente fixadas foi reajustado, sem exclusão de
coletas, em três análises: (i) LOSO restrito às coletas de 18 e 19 de maio;
(ii) treinamento nas coletas de 11 a 13 de maio e teste nas de 18 a 19;
e (iii) treinamento nas de 18 a 19 e teste nas de 11 a 13. Para cada
separação, avaliaram-se as nove entradas, apenas os seis MQ e apenas as
três variáveis ambientais, com os mesmos hiperparâmetros do ET da base
completa. As separações entre períodos testam transferência temporal,
mas incluem vasos que podem aparecer tanto no treinamento quanto no
teste; não medem generalização para vasos novos.

Como sensibilidade à deriva entre dias, repetiu-se o LOSO nas 36 coletas
usando somente os seis MQ após padronização separada por dia. Para cada
sensor $s$ e dia $d$, a média $\mu_{s,d}^{(-v)}$ e o desvio padrão
$\sigma_{s,d}^{(-v)}$ foram estimados apenas nas leituras dos nove vasos
de treinamento da rodada $v$. A transformação das leituras de treino e
teste desse dia seguiu a Equação~\eqref{eq:padronizacao_dia_mq}:
\begin{equation}
    \widetilde{x}_{s,d,i}^{(-v)}=
    \frac{x_{s,d,i}-\mu_{s,d}^{(-v)}}
         {\sigma_{s,d}^{(-v)}},
    \label{eq:padronizacao_dia_mq}
\end{equation}
em que $i$ identifica a leitura. Valores do vaso reservado não
participaram do cálculo da média ou do desvio padrão; desvios padrão
nulos seriam substituídos por 1. Essa transformação exige que o dia de
teste também esteja representado nos demais vasos do treinamento.

Para verificar a influência do recorte proporcional de 70\% da série,
repetiu-se o LOSO apenas com as leituras cujo tempo decorrido estava
entre 200 e 800~s, depois da seleção computacional pela pressão. Esse intervalo é um
recorte fixo de sensibilidade: como não há marca da abertura da válvula
em cada série analisada, não se pode afirmar que contenha exclusivamente
a fase de exposição ao \textit{headspace} em todas as coletas. Foram mantidas as 36
coletas e os dez vasos, e o ET foi reajustado nas três configurações
fixas de entrada descritas acima. O desempenho foi novamente calculado
por leitura e por coleta; a duração efetivamente retida em cada coleta
foi resumida por classe.

Por fim, compararam-se as entradas fixas de nove variáveis com as três
ambientais nas mesmas previsões fora de vaso da base completa. Para cada
vaso $v$, calculou-se $d_v$, a diferença entre as acurácias das duas
configurações naquele vaso. A enumeração bicaudal de todas as $2^{10}$
inversões de sinal foi definida pela Equação~\eqref{eq:inversao_sinais_vaso}:
\begin{equation}
    p_{\mathrm{inv}}=2^{-10}
    \sum_{\boldsymbol{\epsilon}\in\{-1,+1\}^{10}}
    \mathbf{1}\!\left\{
    \left|\sum_{v=1}^{10}\epsilon_v d_v\right|
    \geq
    \left|\sum_{v=1}^{10}d_v\right|
    \right\}.
    \label{eq:inversao_sinais_vaso}
\end{equation}
As diferenças foram calculadas separadamente para acurácia por leitura
dentro de cada vaso e para acurácia por coleta dentro de cada vaso.
Esse teste é exploratório: a interpretação como teste exato exige
simetria das diferenças sob a hipótese nula, e as rodadas LOSO compartilham
vasos no treinamento. Não foi realizado teste de permutação dos rótulos
dos vasos com novo ajuste do classificador em cada alocação.

